%% ma: L12-B06-0002
%% lop: 12
%% chuong: 2
%% bai: 6
%% dang: 12.6.2
%% loai: TN4
%% muc: TH
%% dap_an: "C"
%% nguon: "(NBV)-ĐỀ SỐ 1-MỖI NGÀY 1 ĐỀ THI 2025.pdf (D01, MỖI NGÀY 1 ĐỀ - Đề số 1), Phần I Câu 3"
%% kien_thuc: "Bài 6 (các phép toán vectơ trong không gian, quy tắc hình hộp, độ dài vectơ)"
%% ghi_chu: "Đề gốc có hình lập phương; bản này không cần hình vì tên đỉnh đủ xác định. kho.py trung: 0,56 so với L12-B06-0001 (quy tắc hình hộp); câu này xét 4 khẳng định có cả độ dài vectơ nên coi là khác dạng"
%% trang_thai: cho-duyet
%% ly_do: "Dạng mới đề xuất 12.6.2 (xét khẳng định đúng/sai về vectơ và độ dài vectơ trong hình lập phương) chờ giáo viên duyệt"
\begin{ex}
Cho hình lập phương $ABCD.A'B'C'D'$ cạnh $a$. Khẳng định nào sau đây \textbf{SAI}?
\choice{$\left|\overrightarrow{BD}\right|=a\sqrt2$}{$\left|\overrightarrow{BD'}\right|=a\sqrt3$}{\True $\overrightarrow{AC}+\overrightarrow{A'C'}=\overrightarrow0$}{$\overrightarrow{BA}+\overrightarrow{BC}+\overrightarrow{BB'}=\overrightarrow{BD'}$}
\loigiai{$BD$ là đường chéo của hình vuông cạnh $a$ nên $BD=a\sqrt2$; $BD'$ là đường chéo của hình lập phương nên $BD'=a\sqrt3$; $\overrightarrow{BA}+\overrightarrow{BC}+\overrightarrow{BB'}=\overrightarrow{BD'}$ theo quy tắc hình hộp.
Vì $\overrightarrow{A'C'}=\overrightarrow{AC}$ nên $\overrightarrow{AC}+\overrightarrow{A'C'}=2\overrightarrow{AC}\ne\overrightarrow0$. Khẳng định C sai.}
\end{ex}
